\section{Introduction}

A multiple-choice answer hides more than one decision. A model can first narrow the field to a few answers that still seem possible and only later settle on one final winner. Standard evaluation only shows the last step. It tells us what answer came out, but not when that answer first became a serious option inside the network.

This paper asks a simple question. Does the eventual answer become one of the plausible answers before the model writes the margin that finally separates it from the runner-up? Figure~\ref{fig:flagship} previews the answer. In the pooled primary analysis, by 20\% of normalized depth the eventual answer is already in the plausible answer set for \FlagshipTwentyPlausiblePct{} of evaluable cases, but only \FlagshipTwentyCommittedPct{} are already stably committed.

We answer that question in a very literal way. We build a normal multiple-choice prompt, end it with the suffix \texttt{"Final answer: "}, and then score the exact strings the model could write next. On the main path, that string is the full option text itself, such as \texttt{"photosynthesis"}. So instead of asking an external probe to guess the answer, we watch how the model's own answer strings rise and fall across depth.

Our setting is intentionally narrow and easy to audit. We evaluate nine instruction-tuned open-weight models from the Qwen 2.5 Instruct, Gemma 2 IT, and Llama Instruct families on a 1{,}000-example slice spanning ARC-Challenge, CommonsenseQA, and MMLU abstract algebra. The claim of the paper is limited to that setting. Within it, the pattern is consistent: the eventual answer usually enters the plausible answer set well before it becomes stably isolated from the runner-up. All \ArcCsqaPositiveCells{} evaluable ARC-Challenge and CommonsenseQA model-dataset cells are positive under the preregistered primary rule (Figure~\ref{fig:forest}, Table~\ref{tab:primary-summary}). The MMLU abstract-algebra slice is much smaller, so it is too small to carry much weight on its own.

The paper makes three contributions. First, it gives a simple measurement framework for separating ``still possible'' from ``finally chosen'' without leaving the answer strings that appear in the task itself. Second, it shows a large and consistent gap between those two stages in the pooled curves and in every evaluable ARC-Challenge and CommonsenseQA cell. Third, it shows that small changes to answer format, such as switching from full option text to an answer letter, can change both coverage and winners. That makes those alternatives useful controls, but not replacements for exact option-text scoring.

\begin{figure*}[t]
\centering
\includegraphics[width=\textwidth]{figure1_flagship_plausibility_commitment.pdf}
\caption{\textbf{The eventual answer becomes plausible much earlier than it becomes committed.} The top panel pools all evaluable model-prompt pairs from the primary analysis path (\texttt{semantic\_exact} with identity ordering). The light curve asks: by this depth, in how many cases is the eventual answer already one of the plausible answers? The dark curve asks: by this depth, in how many cases has the model already opened a stable winner-runner margin that it keeps through the end of the network? The band between the curves is the main object in the paper: cases where the answer already looks plausible but the model has not fully committed yet. The interval strip shows a sample of individual cases on the same normalized-depth axis, and the three bottom panels show the same pattern separately for ARC-Challenge, CommonsenseQA, and MMLU abstract algebra.}
\label{fig:flagship}
\end{figure*}
